\documentclass[aspectratio=169, 14pt]{beamer}
\usepackage{stampfly_slides}

\lesson{8}{システムモデリング}{System Modeling}

\begin{document}

% ------------------------------------------------------------------
\begin{frame}
  \titlepage
\end{frame}

% ------------------------------------------------------------------
\begin{frame}{今日のゴール / Today's Goal}
  \begin{block}{目標}
    プラント伝達関数を導出し、モデルに基づいて P ゲインを設計する
  \end{block}

  \vspace{1em}

  \begin{enumerate}
    \item モータ+機体のプラント伝達関数 $G_p(s)$ を理解
    \item P 制御の閉ループ特性（$\omega_n$, $\zeta$）を導出
    \item $\zeta$ から $K_p$ を設計 --- L05/L07 の経験をモデルで裏付け
  \end{enumerate}
\end{frame}

% ------------------------------------------------------------------
\begin{frame}{プラントモデル / Plant Model}
  \begin{center}
    \vspace{-2mm}
    \resizebox{0.85\textwidth}{!}{\input{../tikz/plant_model}}
  \end{center}

  \vspace{-0.5em}

  \begin{exampleblock}{2つのブロック}
    \begin{itemize}
      \item \textbf{Mixer + Motor}: duty 入力 $\to$ トルク
            --- 1次遅れ $\tau_m \approx 20$\,ms
      \item \textbf{Body}: トルク $\to$ 角速度
            --- 剛体回転の積分器 $1/(I \cdot s)$
    \end{itemize}
  \end{exampleblock}
\end{frame}

% ------------------------------------------------------------------
\begin{frame}{$\tau_m$ と $K_m$ の導出 / Deriving $\tau_m$ and $K_m$}
  \vspace{-2mm}
  \begin{columns}[t]
    \begin{column}{0.48\textwidth}
      \begin{block}{モータ時定数 $\tau_m$}\small
        DC モータの電気$\cdot$機械モデルから:
        \vspace{-0.5em}
        \[
          \tau_m = \frac{J_{mp} \cdot R_m}{K_e^2 + D_m R_m}
        \]
        \vspace{-1.5em}
      \end{block}
      \vspace{-0.2em}
      {\scriptsize
      \begin{tabular}{lll}
        \hline
        $J_{mp}$ & 2.01e-8\,kg$\cdot$m$^2$ & 回転子慣性 \\
        $R_m$    & 0.34\,$\Omega$    & 巻線抵抗 \\
        $K_e$    & 6.13e-4\,V/(rad/s) & 逆起電力定数 \\
        $D_m$    & 3.69e-8 & 粘性摩擦 \\
        \hline
        $\tau_m$ & \textbf{0.018\,s} $\approx$ \textbf{0.02\,s} & \\
        \hline
      \end{tabular}}
    \end{column}
    \begin{column}{0.52\textwidth}
      \begin{block}{実効トルクゲイン $K_m$}\small
        $K_f$: ホバー点で duty を $\Delta d$ 変えたときの
        1モータ推力変化量\vspace{-0.3em}
        \[
          K_f = \left.\frac{dF}{dd}\right|_{\!hover}\!,\quad
          K_m = \begin{cases}
            4L \cdot K_f & \text{Roll/Pitch}\\
            4\kappa \cdot K_f & \text{Yaw}
          \end{cases}
        \]
        \vspace{-1.5em}
      \end{block}
      \vspace{-0.2em}
      {\scriptsize
      \begin{tabular}{lll}
        \hline
        $K_f$       & 0.010\,N/duty & 同定値 \\
        $L$         & 0.023\,m & アーム長 \\
        $\kappa$    & 9.71e-3\,m & トルク推力比 \\
        \hline
        $K_{m,rp}$  & \textbf{9.3e-4}\,N$\cdot$m & Roll/Pitch \\
        $K_{m,yaw}$ & \textbf{3.9e-4}\,N$\cdot$m & Yaw \\
        \hline
      \end{tabular}}
    \end{column}
  \end{columns}

  \vspace{0.5em}
  \centering
  実効プラントゲイン: \; $K = K_m / I$
  \quad ($K_{roll}\!=\!102$,\; $K_{pitch}\!=\!70$,\; $K_{yaw}\!=\!19$\,rad/s$^2$)
\end{frame}

% ------------------------------------------------------------------
\begin{frame}{パラメータ一覧 / Parameter Summary}
  \begin{table}
    \centering
    \begin{tabular}{llccc}
      \hline
      \textbf{パラメータ} & \textbf{記号} & \textbf{Roll} & \textbf{Pitch} & \textbf{Yaw} \\
      \hline
      慣性モーメント & $I$ & 9.16e-6 & 13.3e-6 & 20.4e-6\,kg\,$\cdot$\,m$^2$ \\
      トルクゲイン & $K_m$ & 9.3e-4 & 9.3e-4 & 3.9e-4\,N$\cdot$m \\
      モータ時定数 & $\tau_m$ & \multicolumn{3}{c}{0.02\,s（共通）} \\
      \hline
      \textbf{プラントゲイン} & $\bm{K = K_m/I}$
        & \textbf{102} & \textbf{70} & \textbf{19}\,rad/s$^2$ \\
      \hline
    \end{tabular}
  \end{table}

  \vspace{0.5em}

  \begin{exampleblock}{なぜ軸ごとに $K$ が違う？}
    \begin{itemize}
      \item Roll/Pitch: 同じ $K_m$ だが $I_{yy} > I_{xx}$ → Pitch の方が遅い
      \item Yaw: $\kappa \ll L$ なので $K_m$ 自体が小さい → さらに遅い
    \end{itemize}
  \end{exampleblock}
\end{frame}

% ------------------------------------------------------------------
\begin{frame}{P 制御の閉ループ / P Control Closed Loop}
  \begin{center}
    \vspace{-2mm}
    \resizebox{0.85\textwidth}{!}{\input{../tikz/closed_loop_model}}
  \end{center}

  \vspace{-0.5em}

  \begin{block}{閉ループ伝達関数}
    $G_{cl}(s) = \dfrac{K_p K}{\tau_m s^2 + s + K_p K}$
    \hfill $\longleftarrow$ \textbf{2 次系！}
  \end{block}
\end{frame}

% ------------------------------------------------------------------
\begin{frame}{2次系の特性 / Second-Order Characteristics}
  \vspace{-2mm}
  \begin{columns}[c]
    \begin{column}{0.42\textwidth}
      \centering
      \resizebox{\columnwidth}{!}{\input{../tikz/step_response_zeta}}
    \end{column}
    \begin{column}{0.58\textwidth}
      \begin{block}{固有振動数と減衰比}
        \vspace{-0.8em}
        \begin{align*}
          \omega_n &= \sqrt{\frac{K_p K}{\tau_m}}, \quad
          \zeta = \frac{1}{2\sqrt{K_p K \tau_m}}
        \end{align*}
        \vspace{-1.2em}
      \end{block}

      \vspace{0.2em}

      \begin{block}{設計式: $\zeta$ から $K_p$ を逆算}
        \vspace{-0.5em}
        \[
          K_p = \frac{1}{4\zeta^2 K \tau_m}
        \]
        \vspace{-1em}
      \end{block}

      \vspace{0.2em}

      {\small オーバーシュート量
      $\approx e^{-\pi\zeta/\sqrt{1-\zeta^2}}$
      \quad ($\zeta\!=\!0.5 \to 16\%$,\; $\zeta\!=\!0.7 \to 5\%$)}
    \end{column}
  \end{columns}
\end{frame}

% ------------------------------------------------------------------
\begin{frame}{ゲイン設計表 / Gain Design Table}
  \begin{table}
    \centering
    \begin{tabular}{lcccc}
      \hline
      & \textbf{$K$} & \textbf{$K_p=0.5$ の $\zeta$}
      & \textbf{$\zeta=0.7$ の $K_p$} & \textbf{$\zeta=1.0$ の $K_p$} \\
      \hline
      Roll  & 102 & 0.50 & 0.25 & 0.12 \\
      Pitch &  70 & 0.60 & 0.36 & 0.18 \\
      Yaw   &  19 & 1.15 & 1.34 & 0.66 \\
      \hline
    \end{tabular}
  \end{table}

  \vspace{0.5em}

  \begin{exampleblock}{読み取り}
    \begin{itemize}
      \item L05 の $K_p=0.5$ は Roll で $\zeta=0.50$（やや振動的）
      \item Yaw は $\zeta > 1$（過減衰 = 応答が遅い）--- 軸ごとに
            $K_p$ を変えるべき理由！
    \end{itemize}
  \end{exampleblock}
\end{frame}

% ------------------------------------------------------------------
\begin{frame}[fragile]{実習: モデルベースゲイン設計 / Hands-on}
  \begin{lstlisting}[style=sfcpp, title={user\_code.cpp (excerpt)}]
// Model-based Kp design (target zeta = 0.7)
const float tau_m = 0.02f;
const float K_roll  = 102.0f;
const float K_pitch =  70.0f;
const float K_yaw   =  19.0f;
float zeta = 0.7f;

float Kp_roll  = 1.0f/(4*zeta*zeta*K_roll *tau_m);
float Kp_pitch = 1.0f/(4*zeta*zeta*K_pitch*tau_m);
float Kp_yaw   = 1.0f/(4*zeta*zeta*K_yaw  *tau_m);

// Apply per-axis Kp in control loop
float roll_cmd  = Kp_roll  * (target_roll  - ws::gyro_x());
float pitch_cmd = Kp_pitch * (target_pitch - ws::gyro_y());
float yaw_cmd   = Kp_yaw   * (target_yaw   - ws::gyro_z());
  \end{lstlisting}
\end{frame}

% ------------------------------------------------------------------
\begin{frame}{チェックポイント / Checkpoint}
  \begin{alertblock}{確認事項}
    \begin{itemize}
      \item[$\square$] $G_p(s) = K / (s(\tau_m s + 1))$ を説明できる
      \item[$\square$] モデルベース $K_p$ と L05 の $K_p=0.5$ を比較飛行
      \item[$\square$] Roll/Pitch/Yaw の $\zeta$ 差を体感で理解
    \end{itemize}
  \end{alertblock}

  \vspace{1em}

  \begin{block}{次のレッスン}
    Lesson 9: 姿勢推定 --- センサフュージョンで角度を推定
  \end{block}
\end{frame}

\end{document}
